\kafkasection{Hydrogen Molecules can exist in ortho and para states}
\kafkasubsection{The two electrons of $H_2$ in para-hydrogen form a singlet state. The orbital angular momentum can thus only take even values; that is the Hamiltonian for the rotational modes is
\begin{equation}
    \Haml_p = \frac{\hbar^2}{2I}l(l+1),
\end{equation}
where $l = 0,2,4...$ calculate the rotational partition function of para-hydrogen, and evaluate its low and high temperature limits.
}
\kafkasubsection{In ortho-hydrogen the electrons are in a triply degenerate symmetric state, henc
\begin{equation}
    \Haml_o = \frac{\hbar^2}{2I}l(l+1)
\end{equation}
where $l = 1,3,5...$ calculate the rotational partition function of ortho-hydrogen, and evaluate its low and high temperature limits. 
}
\kafkasubsection{For an equilibrium gas of $N$ hydrogen molecules calculate the partition function ignoring virbrational degrees of freedom. Sum over contributions from mixtures of $N_p$ para and $N_o = N - N_p$ ortho-hydrogen particles.}
\kafkasubsection{Write down the expression for the rotational contributions to the internal energy $E_{rot}$, and comment on its low and high temperature limits.}